% --- Archon physics formalization source begin ---
% archon:physics
% archon:covers PhyXMiniProblems/problem_phyx_mini_0242.lean
% archon:source-report reports/phyx_mini/problem_phyx_mini_0242.source.json

\chapter{Physics problem phyx\_mini\_0242}
\label{ch:PhyXMiniProblems_problem_phyx_mini_0242}

\paragraph{Problem source.}
\#\# Physical scenario

A light string with a mass per unit length of \textbackslash{}( 8.00\textbackslash{},\textbackslash{}mathrm\{g/m\} \textbackslash{}) has its ends tied to two walls separated by a distance equal to three-fourths the length of the string (figure). An object of mass \textbackslash{}( m \textbackslash{}) is suspended from the center of the string, putting a tension in the string.The wave speed is to be \textbackslash{}( 60.0\textbackslash{},\textbackslash{}mathrm\{m/s\} \textbackslash{}).

\#\# Question

What should be the mass of the object suspended from the string?

\#\# Answer choices

- A: A: 3.19kg

- B: B: 4.19kg

- C: C: 3.99kg

- D: D: 3.89kg

\#\# Figure caption (auxiliary; use the image as primary evidence)

The image depicts a physics diagram with two vertical walls and a mass suspended between them. A rope is tied to each wall and the mass forms an inverted "V" shape. The left side of the rope is labeled with a length of \textbackslash{}( \textbackslash{}frac\{L\}\{2\} \textbackslash{}), and the same label appears on the right side. The horizontal distance between the two points where the ropes are attached to the walls is labeled \textbackslash{}( \textbackslash{}frac\{3L\}\{4\} \textbackslash{}).

The mass, denoted by \textbackslash{}( m \textbackslash{}), is hanging at the lowest point where the two ropes meet.

\#\# Dataset metadata

Wave Properties; Physical Model Grounding Reasoning, Multi-Formula Reasoning

\paragraph{Recorded answer/context.}
D: D: 3.89kg

\paragraph{Figure/image path.}
/root/proposal\_for\_physic/science-mango/phyx\_mini\_run/phyx\_data/test\_image/242.png

\paragraph{Formalization target.}
create a compiling Lean file with sorry bodies at `PhyXMiniProblems/problem\_phyx\_mini\_0242.lean`. The Lean declarations must preserve the physical quantities, dimensions or dimensional roles, figure labels, governing-law hypotheses, and final relation expressed by this problem.
Use Mathlib/Physlib names found through LeanExplore where available. If a physics API is missing, introduce faithful local abstractions rather than scalar placeholder aliases.

\begin{theorem}[Physics formalization target]
\label{thm:physics:phyx_mini_0242:target}
\lean{PhyXMiniProblems.ProblemPhyXMini0242.problem_phyx_mini_0242}
\uses{def:physics:phyx-mini-0242:phyxminiproblems-problemphyxmini0242-massinkilograms, def:physics:phyx-mini-0242:phyxminiproblems-problemphyxmini0242-suspendedstringsetup, def:physics:phyx-mini-0242:phyxminiproblems-problemphyxmini0242-matchesproblemandprimaryfigure, def:physics:phyx-mini-0242:phyxminiproblems-problemphyxmini0242-satisfiesstraightsegmentgeometry, def:physics:phyx-mini-0242:phyxminiproblems-problemphyxmini0242-usesstandardgravity, def:physics:phyx-mini-0242:phyxminiproblems-problemphyxmini0242-haspositivephysicalparameters, def:physics:phyx-mini-0242:phyxminiproblems-problemphyxmini0242-satisfieswaveandstaticequilibriumlaws, def:physics:phyx-mini-0242:phyxminiproblems-problemphyxmini0242-recordedanswerchoice, def:physics:phyx-mini-0242:phyxminiproblems-problemphyxmini0242-matchesanswerchoice}
The assigned autoformalize agent should translate this physics problem into Lean declarations in the covered file, with theorem and lemma proof bodies written as `by sorry`.
\end{theorem}
\begin{proof}
This is an autoformalization task, not a proof task. Produce statements that can later be proved by the physics prover without weakening the physical model.
\end{proof}

% --- Archon declaration topology begin ---
\section{Declaration topology}
\begin{definition}[Length Quantity]
  \label{def:physics:phyx-mini-0242:phyxminiproblems-problemphyxmini0242-lengthquantity}
  \lean{PhyXMiniProblems.ProblemPhyXMini0242.LengthQuantity}
  A physical length
\end{definition}
\begin{proof}
  This is the declared model definition of Length Quantity.
\end{proof}
\begin{definition}[Mass Quantity]
  \label{def:physics:phyx-mini-0242:phyxminiproblems-problemphyxmini0242-massquantity}
  \lean{PhyXMiniProblems.ProblemPhyXMini0242.MassQuantity}
  A physical mass
\end{definition}
\begin{proof}
  This is the declared model definition of Mass Quantity.
\end{proof}
\begin{definition}[Linear Mass Density Quantity]
  \label{def:physics:phyx-mini-0242:phyxminiproblems-problemphyxmini0242-linearmassdensityquantity}
  \lean{PhyXMiniProblems.ProblemPhyXMini0242.LinearMassDensityQuantity}
  A physical mass per unit length, with dimension M L⁻¹
\end{definition}
\begin{proof}
  This is the declared model definition of Linear Mass Density Quantity.
\end{proof}
\begin{definition}[Speed Quantity]
  \label{def:physics:phyx-mini-0242:phyxminiproblems-problemphyxmini0242-speedquantity}
  \lean{PhyXMiniProblems.ProblemPhyXMini0242.SpeedQuantity}
  A physical propagation speed, with dimension L T⁻¹
\end{definition}
\begin{proof}
  This is the declared model definition of Speed Quantity.
\end{proof}
\begin{definition}[Acceleration Quantity]
  \label{def:physics:phyx-mini-0242:phyxminiproblems-problemphyxmini0242-accelerationquantity}
  \lean{PhyXMiniProblems.ProblemPhyXMini0242.AccelerationQuantity}
  A physical acceleration magnitude, with dimension L T⁻²
\end{definition}
\begin{proof}
  This is the declared model definition of Acceleration Quantity.
\end{proof}
\begin{definition}[Force Quantity]
  \label{def:physics:phyx-mini-0242:phyxminiproblems-problemphyxmini0242-forcequantity}
  \lean{PhyXMiniProblems.ProblemPhyXMini0242.ForceQuantity}
  A physical force magnitude, used for tension and weight
\end{definition}
\begin{proof}
  This is the declared model definition of Force Quantity.
\end{proof}
\begin{definition}[quantity Readout]
  \label{def:physics:phyx-mini-0242:phyxminiproblems-problemphyxmini0242-quantityreadout}
  \lean{PhyXMiniProblems.ProblemPhyXMini0242.quantityReadout}
  Read a nonnegative dimensionful quantity in a chosen coherent unit system
\end{definition}
\begin{proof}
  This is the declared model definition of quantity Readout.
\end{proof}
\begin{definition}[length In Meters]
  \label{def:physics:phyx-mini-0242:phyxminiproblems-problemphyxmini0242-lengthinmeters}
  \lean{PhyXMiniProblems.ProblemPhyXMini0242.lengthInMeters}
  \uses{def:physics:phyx-mini-0242:phyxminiproblems-problemphyxmini0242-lengthquantity, def:physics:phyx-mini-0242:phyxminiproblems-problemphyxmini0242-quantityreadout}
  Meter readout of a physical length
\end{definition}
\begin{proof}
  This is the declared model definition of length In Meters.
\end{proof}
\begin{definition}[mass In Kilograms]
  \label{def:physics:phyx-mini-0242:phyxminiproblems-problemphyxmini0242-massinkilograms}
  \lean{PhyXMiniProblems.ProblemPhyXMini0242.massInKilograms}
  \uses{def:physics:phyx-mini-0242:phyxminiproblems-problemphyxmini0242-massquantity, def:physics:phyx-mini-0242:phyxminiproblems-problemphyxmini0242-quantityreadout}
  Kilogram readout of a physical mass
\end{definition}
\begin{proof}
  This is the declared model definition of mass In Kilograms.
\end{proof}
\begin{definition}[linear Mass Density In Kilograms Per Meter]
  \label{def:physics:phyx-mini-0242:phyxminiproblems-problemphyxmini0242-linearmassdensityinkilogramspermeter}
  \lean{PhyXMiniProblems.ProblemPhyXMini0242.linearMassDensityInKilogramsPerMeter}
  \uses{def:physics:phyx-mini-0242:phyxminiproblems-problemphyxmini0242-linearmassdensityquantity, def:physics:phyx-mini-0242:phyxminiproblems-problemphyxmini0242-quantityreadout}
  Kilograms-per-meter readout of a physical linear mass density
\end{definition}
\begin{proof}
  This is the declared model definition of linear Mass Density In Kilograms Per Meter.
\end{proof}
\begin{definition}[speed In Meters Per Second]
  \label{def:physics:phyx-mini-0242:phyxminiproblems-problemphyxmini0242-speedinmeterspersecond}
  \lean{PhyXMiniProblems.ProblemPhyXMini0242.speedInMetersPerSecond}
  \uses{def:physics:phyx-mini-0242:phyxminiproblems-problemphyxmini0242-speedquantity, def:physics:phyx-mini-0242:phyxminiproblems-problemphyxmini0242-quantityreadout}
  Meters-per-second readout of a physical speed
\end{definition}
\begin{proof}
  This is the declared model definition of speed In Meters Per Second.
\end{proof}
\begin{definition}[acceleration In Meters Per Second Squared]
  \label{def:physics:phyx-mini-0242:phyxminiproblems-problemphyxmini0242-accelerationinmeterspersecondsquared}
  \lean{PhyXMiniProblems.ProblemPhyXMini0242.accelerationInMetersPerSecondSquared}
  \uses{def:physics:phyx-mini-0242:phyxminiproblems-problemphyxmini0242-accelerationquantity, def:physics:phyx-mini-0242:phyxminiproblems-problemphyxmini0242-quantityreadout}
  Meters-per-second-squared readout of an acceleration magnitude
\end{definition}
\begin{proof}
  This is the declared model definition of acceleration In Meters Per Second Squared.
\end{proof}
\begin{definition}[force In Newtons]
  \label{def:physics:phyx-mini-0242:phyxminiproblems-problemphyxmini0242-forceinnewtons}
  \lean{PhyXMiniProblems.ProblemPhyXMini0242.forceInNewtons}
  \uses{def:physics:phyx-mini-0242:phyxminiproblems-problemphyxmini0242-forcequantity, def:physics:phyx-mini-0242:phyxminiproblems-problemphyxmini0242-quantityreadout}
  Newton readout of a force magnitude in coherent SI base units
\end{definition}
\begin{proof}
  This is the declared model definition of force In Newtons.
\end{proof}
\begin{definition}[String Side]
  \label{def:physics:phyx-mini-0242:phyxminiproblems-problemphyxmini0242-stringside}
  \lean{PhyXMiniProblems.ProblemPhyXMini0242.StringSide}
  The two straight halves of the string shown in the figure
\end{definition}
\begin{proof}
  This is the declared model definition of String Side.
\end{proof}
\begin{definition}[Figure Point]
  \label{def:physics:phyx-mini-0242:phyxminiproblems-problemphyxmini0242-figurepoint}
  \lean{PhyXMiniProblems.ProblemPhyXMini0242.FigurePoint}
  The three labeled geometric points visible in the primary figure
\end{definition}
\begin{proof}
  This is the declared model definition of Figure Point.
\end{proof}
\begin{definition}[wall Tie Point]
  \label{def:physics:phyx-mini-0242:phyxminiproblems-problemphyxmini0242-stringside-walltiepoint}
  \lean{PhyXMiniProblems.ProblemPhyXMini0242.StringSide.wallTiePoint}
  \uses{def:physics:phyx-mini-0242:phyxminiproblems-problemphyxmini0242-stringside, def:physics:phyx-mini-0242:phyxminiproblems-problemphyxmini0242-figurepoint}
  The wall tie point belonging to a given half of the string
\end{definition}
\begin{proof}
  This is the declared model definition of wall Tie Point.
\end{proof}
\begin{definition}[Suspended String Setup]
  \label{def:physics:phyx-mini-0242:phyxminiproblems-problemphyxmini0242-suspendedstringsetup}
  \lean{PhyXMiniProblems.ProblemPhyXMini0242.SuspendedStringSetup}
  \uses{def:physics:phyx-mini-0242:phyxminiproblems-problemphyxmini0242-lengthquantity, def:physics:phyx-mini-0242:phyxminiproblems-problemphyxmini0242-massquantity, def:physics:phyx-mini-0242:phyxminiproblems-problemphyxmini0242-linearmassdensityquantity, def:physics:phyx-mini-0242:phyxminiproblems-problemphyxmini0242-speedquantity, def:physics:phyx-mini-0242:phyxminiproblems-problemphyxmini0242-accelerationquantity, def:physics:phyx-mini-0242:phyxminiproblems-problemphyxmini0242-forcequantity, def:physics:phyx-mini-0242:phyxminiproblems-problemphyxmini0242-stringside, def:physics:phyx-mini-0242:phyxminiproblems-problemphyxmini0242-figurepoint}
  The complete state of the string and suspended object. The planar coordinates are scalar readouts in meters; all lengths and mechanical quantities remain dimensionful objects. No field assigns the suspended mass its requested value
\end{definition}
\begin{proof}
  This is the declared model definition of Suspended String Setup.
\end{proof}
\begin{definition}[x Coordinate In Meters]
  \label{def:physics:phyx-mini-0242:phyxminiproblems-problemphyxmini0242-xcoordinateinmeters}
  \lean{PhyXMiniProblems.ProblemPhyXMini0242.xCoordinateInMeters}
  \uses{def:physics:phyx-mini-0242:phyxminiproblems-problemphyxmini0242-figurepoint, def:physics:phyx-mini-0242:phyxminiproblems-problemphyxmini0242-suspendedstringsetup}
  Horizontal coordinate, in meters, of a labeled figure point
\end{definition}
\begin{proof}
  This is the declared model definition of x Coordinate In Meters.
\end{proof}
\begin{definition}[y Coordinate In Meters]
  \label{def:physics:phyx-mini-0242:phyxminiproblems-problemphyxmini0242-ycoordinateinmeters}
  \lean{PhyXMiniProblems.ProblemPhyXMini0242.yCoordinateInMeters}
  \uses{def:physics:phyx-mini-0242:phyxminiproblems-problemphyxmini0242-figurepoint, def:physics:phyx-mini-0242:phyxminiproblems-problemphyxmini0242-suspendedstringsetup}
  Vertical coordinate, in meters, of a labeled figure point
\end{definition}
\begin{proof}
  This is the declared model definition of y Coordinate In Meters.
\end{proof}
\begin{definition}[Matches Problem And Primary Figure]
  \label{def:physics:phyx-mini-0242:phyxminiproblems-problemphyxmini0242-matchesproblemandprimaryfigure}
  \lean{PhyXMiniProblems.ProblemPhyXMini0242.MatchesProblemAndPrimaryFigure}
  \uses{def:physics:phyx-mini-0242:phyxminiproblems-problemphyxmini0242-lengthinmeters, def:physics:phyx-mini-0242:phyxminiproblems-problemphyxmini0242-linearmassdensityinkilogramspermeter, def:physics:phyx-mini-0242:phyxminiproblems-problemphyxmini0242-speedinmeterspersecond, def:physics:phyx-mini-0242:phyxminiproblems-problemphyxmini0242-stringside, def:physics:phyx-mini-0242:phyxminiproblems-problemphyxmini0242-suspendedstringsetup, def:physics:phyx-mini-0242:phyxminiproblems-problemphyxmini0242-xcoordinateinmeters, def:physics:phyx-mini-0242:phyxminiproblems-problemphyxmini0242-ycoordinateinmeters}
  Data read directly from the prose and primary figure. It records the 8.00 g/m = 1/125 kg/m density, the 60.0 m/s wave speed, the two labels L/2, the wall separation 3L/4, the centered object, and the displayed inverted-V geometry. It contains no mass value
\end{definition}
\begin{proof}
  This is the declared model definition of Matches Problem And Primary Figure.
\end{proof}
\begin{definition}[Satisfies Straight Segment Geometry]
  \label{def:physics:phyx-mini-0242:phyxminiproblems-problemphyxmini0242-satisfiesstraightsegmentgeometry}
  \lean{PhyXMiniProblems.ProblemPhyXMini0242.SatisfiesStraightSegmentGeometry}
  \uses{def:physics:phyx-mini-0242:phyxminiproblems-problemphyxmini0242-lengthinmeters, def:physics:phyx-mini-0242:phyxminiproblems-problemphyxmini0242-stringside, def:physics:phyx-mini-0242:phyxminiproblems-problemphyxmini0242-suspendedstringsetup}
  Each string half is a straight segment, so its length, horizontal projection, and vertical drop obey the planar Pythagorean relation. This is a governing geometric law; it does not assume the derived sqrt 7 / 4 direction ratio
\end{definition}
\begin{proof}
  This is the declared model definition of Satisfies Straight Segment Geometry.
\end{proof}
\begin{definition}[Uses Standard Gravity]
  \label{def:physics:phyx-mini-0242:phyxminiproblems-problemphyxmini0242-usesstandardgravity}
  \lean{PhyXMiniProblems.ProblemPhyXMini0242.UsesStandardGravity}
  \uses{def:physics:phyx-mini-0242:phyxminiproblems-problemphyxmini0242-accelerationinmeterspersecondsquared, def:physics:phyx-mini-0242:phyxminiproblems-problemphyxmini0242-suspendedstringsetup}
  The standard classroom calibration g = 9.80 m/s²
\end{definition}
\begin{proof}
  This is the declared model definition of Uses Standard Gravity.
\end{proof}
\begin{definition}[Has Positive Physical Parameters]
  \label{def:physics:phyx-mini-0242:phyxminiproblems-problemphyxmini0242-haspositivephysicalparameters}
  \lean{PhyXMiniProblems.ProblemPhyXMini0242.HasPositivePhysicalParameters}
  \uses{def:physics:phyx-mini-0242:phyxminiproblems-problemphyxmini0242-lengthinmeters, def:physics:phyx-mini-0242:phyxminiproblems-problemphyxmini0242-massinkilograms, def:physics:phyx-mini-0242:phyxminiproblems-problemphyxmini0242-linearmassdensityinkilogramspermeter, def:physics:phyx-mini-0242:phyxminiproblems-problemphyxmini0242-speedinmeterspersecond, def:physics:phyx-mini-0242:phyxminiproblems-problemphyxmini0242-accelerationinmeterspersecondsquared, def:physics:phyx-mini-0242:phyxminiproblems-problemphyxmini0242-forceinnewtons, def:physics:phyx-mini-0242:phyxminiproblems-problemphyxmini0242-stringside, def:physics:phyx-mini-0242:phyxminiproblems-problemphyxmini0242-suspendedstringsetup}
  Positivity conditions for a nondegenerate taut-string configuration
\end{definition}
\begin{proof}
  This is the declared model definition of Has Positive Physical Parameters.
\end{proof}
\begin{definition}[Satisfies Wave And Static Equilibrium Laws]
  \label{def:physics:phyx-mini-0242:phyxminiproblems-problemphyxmini0242-satisfieswaveandstaticequilibriumlaws}
  \lean{PhyXMiniProblems.ProblemPhyXMini0242.SatisfiesWaveAndStaticEquilibriumLaws}
  \uses{def:physics:phyx-mini-0242:phyxminiproblems-problemphyxmini0242-lengthinmeters, def:physics:phyx-mini-0242:phyxminiproblems-problemphyxmini0242-massinkilograms, def:physics:phyx-mini-0242:phyxminiproblems-problemphyxmini0242-linearmassdensityinkilogramspermeter, def:physics:phyx-mini-0242:phyxminiproblems-problemphyxmini0242-speedinmeterspersecond, def:physics:phyx-mini-0242:phyxminiproblems-problemphyxmini0242-accelerationinmeterspersecondsquared, def:physics:phyx-mini-0242:phyxminiproblems-problemphyxmini0242-forceinnewtons, def:physics:phyx-mini-0242:phyxminiproblems-problemphyxmini0242-stringside, def:physics:phyx-mini-0242:phyxminiproblems-problemphyxmini0242-suspendedstringsetup}
  The taut-string wave law, weight law, light-string tension equality, and vertical force balance. The last equation resolves each tension along its straight string segment; no numerical mass or answer choice occurs here
\end{definition}
\begin{proof}
  This is the declared model definition of Satisfies Wave And Static Equilibrium Laws.
\end{proof}
\begin{lemma}[vertical Component Ratio eq sqrt Seven Over Four]
  \label{lem:physics:phyx-mini-0242:phyxminiproblems-problemphyxmini0242-verticalcomponentratio-eq-sqrtsevenoverfour}
  \lean{PhyXMiniProblems.ProblemPhyXMini0242.verticalComponentRatio_eq_sqrtSevenOverFour}
  \uses{def:physics:phyx-mini-0242:phyxminiproblems-problemphyxmini0242-lengthinmeters, def:physics:phyx-mini-0242:phyxminiproblems-problemphyxmini0242-stringside, def:physics:phyx-mini-0242:phyxminiproblems-problemphyxmini0242-suspendedstringsetup, def:physics:phyx-mini-0242:phyxminiproblems-problemphyxmini0242-matchesproblemandprimaryfigure, def:physics:phyx-mini-0242:phyxminiproblems-problemphyxmini0242-satisfiesstraightsegmentgeometry, def:physics:phyx-mini-0242:phyxminiproblems-problemphyxmini0242-haspositivephysicalparameters}
  The figure labels and straight-segment geometry imply that the vertical direction ratio of either string half is sqrt 7 / 4. This relation is derived rather than included among the figure readouts
\end{lemma}
\begin{proof}
  The conclusion follows from the governing relations and figure readouts named in the statement of vertical Component Ratio eq sqrt Seven Over Four.
\end{proof}
\begin{lemma}[suspended Mass In Kilograms eq exact]
  \label{lem:physics:phyx-mini-0242:phyxminiproblems-problemphyxmini0242-suspendedmassinkilograms-eq-exact}
  \lean{PhyXMiniProblems.ProblemPhyXMini0242.suspendedMassInKilograms_eq_exact}
  \uses{def:physics:phyx-mini-0242:phyxminiproblems-problemphyxmini0242-massinkilograms, def:physics:phyx-mini-0242:phyxminiproblems-problemphyxmini0242-suspendedstringsetup, def:physics:phyx-mini-0242:phyxminiproblems-problemphyxmini0242-matchesproblemandprimaryfigure, def:physics:phyx-mini-0242:phyxminiproblems-problemphyxmini0242-satisfiesstraightsegmentgeometry, def:physics:phyx-mini-0242:phyxminiproblems-problemphyxmini0242-usesstandardgravity, def:physics:phyx-mini-0242:phyxminiproblems-problemphyxmini0242-haspositivephysicalparameters, def:physics:phyx-mini-0242:phyxminiproblems-problemphyxmini0242-satisfieswaveandstaticequilibriumlaws}
  The unrounded mass derived from T = μv² and vertical equilibrium is 72 sqrt(7) / 49 kilograms
\end{lemma}
\begin{proof}
  The conclusion follows from the governing relations and figure readouts named in the statement of suspended Mass In Kilograms eq exact.
\end{proof}
\begin{definition}[Answer Choice]
  \label{def:physics:phyx-mini-0242:phyxminiproblems-problemphyxmini0242-answerchoice}
  \lean{PhyXMiniProblems.ProblemPhyXMini0242.AnswerChoice}
  Labels of the four answers printed with the problem
\end{definition}
\begin{proof}
  This is the declared model definition of Answer Choice.
\end{proof}
\begin{definition}[kilograms]
  \label{def:physics:phyx-mini-0242:phyxminiproblems-problemphyxmini0242-answerchoice-kilograms}
  \lean{PhyXMiniProblems.ProblemPhyXMini0242.AnswerChoice.kilograms}
  \uses{def:physics:phyx-mini-0242:phyxminiproblems-problemphyxmini0242-answerchoice}
  Mass in kilograms printed beside each answer label
\end{definition}
\begin{proof}
  This is the declared model definition of kilograms.
\end{proof}
\begin{definition}[recorded Answer Choice]
  \label{def:physics:phyx-mini-0242:phyxminiproblems-problemphyxmini0242-recordedanswerchoice}
  \lean{PhyXMiniProblems.ProblemPhyXMini0242.recordedAnswerChoice}
  \uses{def:physics:phyx-mini-0242:phyxminiproblems-problemphyxmini0242-answerchoice}
  The answer label recorded in the supplied dataset metadata
\end{definition}
\begin{proof}
  This is the declared model definition of recorded Answer Choice.
\end{proof}
\begin{definition}[Matches Answer Choice]
  \label{def:physics:phyx-mini-0242:phyxminiproblems-problemphyxmini0242-matchesanswerchoice}
  \lean{PhyXMiniProblems.ProblemPhyXMini0242.MatchesAnswerChoice}
  \uses{def:physics:phyx-mini-0242:phyxminiproblems-problemphyxmini0242-massquantity, def:physics:phyx-mini-0242:phyxminiproblems-problemphyxmini0242-massinkilograms, def:physics:phyx-mini-0242:phyxminiproblems-problemphyxmini0242-answerchoice}
  Agreement with a kilogram value displayed to the nearest hundredth. The half-hundredth tolerance avoids identifying an unrounded physical value with the printed decimal exactly
\end{definition}
\begin{proof}
  This is the declared model definition of Matches Answer Choice.
\end{proof}
% --- Archon declaration topology end ---
% --- Archon physics formalization source end ---
